% Claim proof file for row claim_3_18 -- "MarginalizationAndIntervention".
% Auto-generated stub by scaffold/solve_chapter.py at row start. A manager
% agent dedicated to writing the TeX proof fills in the proof body below.
%
% This file is a sibling of `main.tex` inside `<subsection>/tex/`. The
% `subfiles` package lets it render both standalone and as part of
% `main.tex` via `\subfile{...}`.
%
% The statement is restated above the proof so the file is self-contained
% when rendered on its own. The proof must:
%   - Stay close to the lecture notes' proof if one exists (always search
%     the LN first for a `\begin{proof}` block immediately following the
%     claim; copy / lightly edit when present).
%   - Otherwise use the LN paradigm and cite prior claims by their `ref`.
%   - Be self-contained enough that a mathematician can follow it without
%     re-reading the LN.
%   - Have no `sorry`, `omitted`, or "obvious" shortcuts.
%
% Part (i) is adapted nearly verbatim from the LN's own proof block at
% lecture-notes/lecture_notes/graphs.tex:1132-1163, with three light edits:
%   - the `Left:` / `Right:` labels are rewritten as `LHS:` / `RHS:` to
%     match Parts (ii) and (iii);
%   - the LN's bare prose references to "marginalization" and "hard
%     intervention" are pinned to def \ref{def:G_marginalization} and
%     def \ref{def:G_hard_intervention} respectively;
%   - the bidirected-edge paragraph's $\ul{v}, \ol{v} \notin W_1$ clause
%     on the RHS is acknowledged to come from the symmetrised reading of
%     def \ref{def:G_hard_intervention}, item~iv.\ (cf.\ the remark there
%     and the registered deviation noted in claim_3_4's proof).
% Parts (ii) and (iii) are written from scratch in field-by-field LN-paradigm
% structure parallel to Part (i), with the inner operator replaced by
% $\doit(I_{W_1})$ (def \ref{def:cdmg_intervention_nodes}) and $\spl(W_1)$
% (def \ref{def:G_node-splitting}) respectively.

\documentclass[main]{subfiles}

\begin{document}
% ref: claim_3_18 (proof, with statement restated for readability)
% title: MarginalizationAndIntervention
\def\rowref{claim\_3\_18}
\phantomsection\label{claim_3_18}

% --- statement (restated) ---------------------------------------------------
\begin{Lem}[Marginalization and intervention commute]\label{marginalization-and-intervention-commute}
    Let $G = (J, V, E, L)$ be a CDMG (def \ref{def-cdmg}). The lemma comprises three commutativity statements -- one per intervention operator: hard intervention (def \ref{def:G_hard_intervention}), adding intervention nodes (def \ref{def:cdmg_intervention_nodes}), and node-splitting (def \ref{def:G_node-splitting}). Each commutes with marginalization (def \ref{def:G_marginalization}). The disjointness side condition
    \[ W_1 \cap W_2 \;=\; \emptyset \]
    is in force throughout. The typing on $W_1$ varies by part: $W_1 \ins J \cup V$ in parts (i) and (ii) (matching the preconditions of def \ref{def:G_hard_intervention} and def \ref{def:cdmg_intervention_nodes} respectively), and $W_1 \ins V$ in part (iii) (matching the precondition of def \ref{def:G_node-splitting}). In all three parts $W_2 \ins V$ matches the precondition of the marginalization operator $G^{\sm W_2}$ (def \ref{def:G_marginalization}). All quantifiers ``for every CDMG $G$'', ``for every $W_1$'', and ``for every $W_2$'' (subject to the typing and disjointness above) are read \emph{universally}; the degenerate cases $W_1 = \emptyset$, $W_2 = \emptyset$, and $W_1 = W_2 = \emptyset$ are admitted.

    The following three equalities of CDMGs hold:
    \begin{enumerate}
        \item[(i)] \emph{(Hard intervention.)} For every $W_1 \ins J \cup V$ and $W_2 \ins V$ with $W_1 \cap W_2 = \emptyset$:
              \[ \lp G_{\doit(W_1)} \rp^{\sm W_2} = \lp G^{\sm W_2} \rp_{\doit(W_1)}. \]
        \item[(ii)] \emph{(Adding intervention nodes.)} For every $W_1 \ins J \cup V$ and $W_2 \ins V$ with $W_1 \cap W_2 = \emptyset$:
              \[ \lp G_{\doit(I_{W_1})} \rp^{\sm W_2} = \lp G^{\sm W_2} \rp_{\doit(I_{W_1})}. \]
        \item[(iii)] \emph{(Node-splitting.)} For every $W_1 \ins V$ and $W_2 \ins V$ with $W_1 \cap W_2 = \emptyset$:
              \[ \lp G_{\spl(W_1)} \rp^{\sm W_2} = \lp G^{\sm W_2} \rp_{\spl(W_1)}. \]
    \end{enumerate}
    Each equality in (i), (ii), (iii) is understood componentwise on the four components $(J, V, E, L)$ of def \ref{def-cdmg}: equal input-node sets, equal output-node sets, equal directed-edge sets, and equal bidirected-edge sets; the conjunctive unpacking into these four field-by-field equalities is deferred to the proof.
    % This shows that the following notations are unambiguous:
    % \[G(V\sm (W_1 \cup W_2)|\doit(W_1 \cup J))=G_{\doit(W_1)}^{\sm W_2}.\]
\end{Lem}

% --- proof ------------------------------------------------------------------
\begin{proof}
We prove the three parts in turn. Each part establishes the displayed equality of CDMGs (def \ref{def-cdmg}) componentwise on the four components $(J, V, E, L)$, in the field-by-field LN paradigm of def \ref{def:G_marginalization}, def \ref{def:G_hard_intervention}, def \ref{def:cdmg_intervention_nodes}, and def \ref{def:G_node-splitting}.

\medskip\noindent\textbf{Part (i) -- hard intervention.}
We verify that all four components of $\lp G_{\doit(W_1)} \rp^{\sm W_2}$ and $\lp G^{\sm W_2} \rp_{\doit(W_1)}$ agree, with $W_1 \ins J \cup V$, $W_2 \ins V$, and $W_1 \cap W_2 = \emptyset$.

\smallskip\noindent\emph{Well-typedness.}
The inner operator on the LHS is $\doit(W_1)$ (def \ref{def:G_hard_intervention}), applied to $G$ with the precondition $W_1 \ins J \cup V$; the resulting CDMG $G_{\doit(W_1)}$ has output-node set $V \sm W_1$ (item~ii.\ of def \ref{def:G_hard_intervention}). The outer marginalization $(\cdot)^{\sm W_2}$ requires $W_2$ to be a subset of the output-node set; this holds because $W_2 \ins V$ and $W_1 \cap W_2 = \emptyset$ give $W_2 \ins V \sm W_1$. Symmetrically on the RHS: $G^{\sm W_2}$ has carrier $J \cup (V \sm W_2)$ (def \ref{def:G_marginalization}, items~i.--ii.); the outer $\doit(W_1)$ requires $W_1 \ins J \cup (V \sm W_2)$, which holds because $W_1 \ins J \cup V$ and $W_1 \cap W_2 = \emptyset$ give $W_1 \ins J \cup (V \sm W_2)$.

\medskip\noindent\emph{Input nodes.}
Both sides have $J \cup W_1$: marginalization does not change input nodes (def \ref{def:G_marginalization}, item~i., gives $J^{\sm W_2} = J$), and hard intervention adds $W_1$ to them (def \ref{def:G_hard_intervention}, item~i., gives $J_{\doit(W_1)} = J \cup W_1$). Concretely,
\[
    J_{\lp G_{\doit(W_1)} \rp^{\sm W_2}}
        \;=\; J_{G_{\doit(W_1)}}
        \;=\; J \cup W_1
        \;=\; J^{\sm W_2} \cup W_1
        \;=\; J_{\lp G^{\sm W_2} \rp_{\doit(W_1)}}. \quad \checkmark
\]

\medskip\noindent\emph{Output nodes.}
LHS: $V_{G_{\doit(W_1)}} \sm W_2 = (V \sm W_1) \sm W_2 = V \sm (W_1 \cup W_2)$. RHS: $V^{\sm W_2} \sm W_1 = (V \sm W_2) \sm W_1 = V \sm (W_1 \cup W_2)$. Both equalities use def \ref{def:G_marginalization}, item~ii., and def \ref{def:G_hard_intervention}, item~ii., together with the standard set identity $(A \sm B) \sm C = A \sm (B \cup C)$. \checkmark

\medskip\noindent\emph{Directed edges.}
Let $\ul{v}, \ol{v} \in J \cup W_1 \cup (V \sm (W_1 \cup W_2))$ with $\ol{v} \in V \sm (W_1 \cup W_2)$ (the head of any directed edge in either side must lie in the common output-node set).

LHS: $\ul{v} \tuh \ol{v} \in (E_{\doit(W_1)})^{\sm W_2}$ iff (by def \ref{def:G_marginalization}, item~iii.) there is a directed walk
\[ \ul{v} \tuh a_1 \tuh \cdots \tuh a_{n-1} \tuh \ol{v} \]
in $G_{\doit(W_1)}$ with $n \ge 1$ and all intermediate nodes $a_1, \dots, a_{n-1} \in W_2$.
By def \ref{def:G_hard_intervention}, item~iii., $E_{\doit(W_1)} = E \sm \lC (v_1, v_2) \in E \st v_2 \in W_1 \rC$, i.e.\ $E_{\doit(W_1)}$ consists of those edges of $E$ whose head does \emph{not} lie in $W_1$. Since every consecutive edge along the walk has its head in $\lC a_1, \ldots, a_{n-1}, \ol{v} \rC \ins W_2 \cup (V \sm (W_1 \cup W_2)) \ins V \sm W_1$ (using $W_1 \cap W_2 = \emptyset$ for the intermediate nodes and $\ol{v} \notin W_1$ by typing for the final node), every such edge of $E_{\doit(W_1)}$ is also an edge of $E$. Hence the walk is also a walk in $G$ through $W_2$. Conversely, any walk in $G$ from $\ul{v}$ to $\ol{v}$ with intermediate nodes in $W_2$ has every edge's head in $W_2 \cup \lC \ol{v} \rC \ins V \sm W_1$, so each such edge survives the removal of edges with head in $W_1$ and lies in $E_{\doit(W_1)}$; thus the walk lifts to a walk in $G_{\doit(W_1)}$ through $W_2$. So the LHS condition holds iff there exists a directed walk in $G$ from $\ul{v}$ to $\ol{v}$ with intermediate nodes in $W_2$.

RHS: $\ul{v} \tuh \ol{v} \in (E^{\sm W_2})_{\doit(W_1)}$ iff (by def \ref{def:G_hard_intervention}, item~iii., applied to $G^{\sm W_2}$) $\ul{v} \tuh \ol{v} \in E^{\sm W_2}$ and $\ol{v} \notin W_1$.
The latter holds by typing since $\ol{v} \in V \sm (W_1 \cup W_2)$.
The former holds (by def \ref{def:G_marginalization}, item~iii.) iff there is a directed walk in $G$ from $\ul{v}$ to $\ol{v}$ with intermediate nodes in $W_2$.

Both sides reduce to ``there exists a directed walk in $G$ from $\ul{v}$ to $\ol{v}$ with intermediate nodes in $W_2$'', so $(E_{\doit(W_1)})^{\sm W_2} = (E^{\sm W_2})_{\doit(W_1)}$. \checkmark

\medskip\noindent\emph{Bidirected edges.}
Let $\ul{v}, \ol{v} \in V \sm (W_1 \cup W_2)$ with $\ul{v} \neq \ol{v}$ (the irreflexivity constraint on bidirected edges, def \ref{def:G_marginalization}, item~iv.).

LHS: $\ul{v} \huh \ol{v} \in (L_{\doit(W_1)})^{\sm W_2}$ iff (by def \ref{def:G_marginalization}, item~iv.) there is a bifurcation in $G_{\doit(W_1)}$ between $\ul{v}$ and $\ol{v}$ with intermediate nodes in $W_2$.
By def \ref{def:G_hard_intervention}, item~iv., read with the symmetrisation needed for the symmetry axiom of def \ref{def-cdmg} (cf.\ the remark following item~iv.\ of def \ref{def:G_hard_intervention} and the registered deviation flagged in \refrow{claim_3_4}'s proof), $L_{\doit(W_1)}$ consists of those ordered pairs $(v_1, v_2) \in L$ with \emph{both} $v_1 \notin W_1$ and $v_2 \notin W_1$. Every node touched by an edge of the bifurcation lies in $\lC \ul{v}, \ol{v}, a_1, \ldots, a_{n-1} \rC \ins W_2 \cup (V \sm (W_1 \cup W_2)) \ins V \sm W_1$ (using $W_1 \cap W_2 = \emptyset$ for the intermediate nodes and the typing of $\ul{v}, \ol{v}$ for the endpoints), so every such edge lies in $L$ as well as $L_{\doit(W_1)}$. Hence the bifurcation is also a bifurcation in $G$ between $\ul{v}$ and $\ol{v}$ with intermediate nodes in $W_2$. Conversely, any bifurcation in $G$ between $\ul{v}$ and $\ol{v}$ with intermediate nodes in $W_2$ has every touched node in $V \sm W_1$, so every edge survives the $W_1$-incident removal and lies in $L_{\doit(W_1)}$ (or in $E_{\doit(W_1)}$ for the directed-arm edges, by the same argument used for directed edges above); hence the bifurcation lifts to $G_{\doit(W_1)}$. So the LHS condition holds iff there exists a bifurcation in $G$ between $\ul{v}$ and $\ol{v}$ with intermediate nodes in $W_2$.

RHS: $\ul{v} \huh \ol{v} \in (L^{\sm W_2})_{\doit(W_1)}$ iff (by def \ref{def:G_hard_intervention}, item~iv., applied to $G^{\sm W_2}$, in its symmetrised reading) $\ul{v} \huh \ol{v} \in L^{\sm W_2}$ and $\ul{v}, \ol{v} \notin W_1$.
The latter holds by typing. The former holds (by def \ref{def:G_marginalization}, item~iv.) iff there is a bifurcation in $G$ between $\ul{v}$ and $\ol{v}$ with intermediate nodes in $W_2$.

Both sides reduce to ``there exists a bifurcation in $G$ between $\ul{v}$ and $\ol{v}$ with intermediate nodes in $W_2$'', so $(L_{\doit(W_1)})^{\sm W_2} = (L^{\sm W_2})_{\doit(W_1)}$. \checkmark

\smallskip
Combining the four field-by-field equalities, $\lp G_{\doit(W_1)} \rp^{\sm W_2} = \lp G^{\sm W_2} \rp_{\doit(W_1)}$ as CDMGs. This establishes Part (i).

\medskip\noindent\textbf{Part (ii) -- adding intervention nodes.}
We verify that all four components of $\lp G_{\doit(I_{W_1})} \rp^{\sm W_2}$ and $\lp G^{\sm W_2} \rp_{\doit(I_{W_1})}$ agree, with $W_1 \ins J \cup V$, $W_2 \ins V$, and $W_1 \cap W_2 = \emptyset$. The construction of $G_{\doit(I_{W_1})}$ is by def \ref{def:cdmg_intervention_nodes}: a fresh intervention symbol $I_w$ is introduced for each $w \in W_1 \sm J$ (type-disjoint from $J \cup V$ and from every other $I_{w'}$), with the notational shorthand $I_j := j$ for $j \in J \cap W_1$ (no new node added on the $J$-branch).

\smallskip\noindent\emph{Well-typedness.}
The inner $\doit(I_{W_1})$ on the LHS is defined per def \ref{def:cdmg_intervention_nodes} from $W_1 \ins J \cup V$. By items~i.\ and ii.\ of def \ref{def:cdmg_intervention_nodes}, $G_{\doit(I_{W_1})}$ has input nodes $J \cup \lC I_w \st w \in W_1 \sm J \rC$ and output nodes $V$. The outer $(\cdot)^{\sm W_2}$ requires $W_2 \ins V$, which is the standing hypothesis. Symmetrically on the RHS: $G^{\sm W_2}$ has carrier $J \cup (V \sm W_2)$ (def \ref{def:G_marginalization}, items~i.--ii.); the outer $\doit(I_{W_1})$ requires $W_1 \ins J \cup (V \sm W_2)$, which holds because $W_1 \ins J \cup V$ and $W_1 \cap W_2 = \emptyset$ give $W_1 \ins J \cup (V \sm W_2)$.

\smallskip\noindent\emph{Carrier identifications used throughout Part (ii).}
The fresh intervention symbols $\lC I_w \st w \in W_1 \sm J \rC$ live on the same carrier on both sides because, on the RHS, the index set ``$W_1 \sm J^{\sm W_2}$'' coincides with ``$W_1 \sm J$'' (since $J^{\sm W_2} = J$ by def \ref{def:G_marginalization}, item~i.). Throughout the proof we identify the fresh-symbol family generated by the LHS construction with the one generated by the RHS construction via the identity on $W_1 \sm J$.

\medskip\noindent\emph{Input nodes.}
LHS: marginalization (def \ref{def:G_marginalization}, item~i.) preserves input nodes; def \ref{def:cdmg_intervention_nodes}, item~i., gives $J_{G_{\doit(I_{W_1})}} = J \cup \lC I_w \st w \in W_1 \sm J \rC$. Hence
\[
    J_{\lp G_{\doit(I_{W_1})} \rp^{\sm W_2}} \;=\; J_{G_{\doit(I_{W_1})}} \;=\; J \cup \lC I_w \st w \in W_1 \sm J \rC.
\]
RHS: by def \ref{def:cdmg_intervention_nodes}, item~i., applied to $G^{\sm W_2}$,
\[
    J_{\lp G^{\sm W_2} \rp_{\doit(I_{W_1})}} \;=\; J^{\sm W_2} \cup \lC I_w \st w \in W_1 \sm J^{\sm W_2} \rC \;=\; J \cup \lC I_w \st w \in W_1 \sm J \rC,
\]
using $J^{\sm W_2} = J$. Equal. \checkmark

\medskip\noindent\emph{Output nodes.}
LHS: $V_{\lp G_{\doit(I_{W_1})} \rp^{\sm W_2}} = V_{G_{\doit(I_{W_1})}} \sm W_2 = V \sm W_2$, using def \ref{def:cdmg_intervention_nodes}, item~ii.\ ($V_{G_{\doit(I_{W_1})}} = V$) and def \ref{def:G_marginalization}, item~ii.
RHS: $V_{\lp G^{\sm W_2} \rp_{\doit(I_{W_1})}} = V^{\sm W_2} = V \sm W_2$, using def \ref{def:cdmg_intervention_nodes}, item~ii.\ (extension preserves $V$), and def \ref{def:G_marginalization}, item~ii. Equal. \checkmark

\medskip\noindent\emph{Directed edges.}
Throughout this paragraph we work over the carrier of $G_{\doit(I_{W_1})}$, namely $J \cup \lC I_w \st w \in W_1 \sm J \rC \cup V$. The subset $W_2 \ins V$ embeds directly into this carrier (no tag-renaming is needed: $V$ sits inside the extended carrier as itself, per def \ref{def:cdmg_intervention_nodes}, item~ii.).

Let $\ul{v}, \ol{v}$ in the appropriate input or output slots (specifically, $\ol{v}$ in the common output-node set $V \sm W_2$; $\ul{v}$ in the common carrier $J \cup \lC I_w \st w \in W_1 \sm J \rC \cup (V \sm W_2)$).

LHS: $\ul{v} \tuh \ol{v} \in (E_{\doit(I_{W_1})})^{\sm W_2}$ iff (by def \ref{def:G_marginalization}, item~iii.) there exists a directed walk
\[ \ul{v} \tuh a_1 \tuh \cdots \tuh a_{n-1} \tuh \ol{v} \]
in $G_{\doit(I_{W_1})}$ with $n \ge 1$ and intermediate nodes $a_1, \dots, a_{n-1} \in W_2$.
By def \ref{def:cdmg_intervention_nodes}, item~iii., $E_{\doit(I_{W_1})} = E \cup \lC (I_w, w) \st w \in W_1 \sm J \rC$. We separate the original edges of $E$ from the fresh edges $\lC (I_w, w) \st w \in W_1 \sm J \rC$.

\emph{Claim: any walk in $G_{\doit(I_{W_1})}$ from $\ul{v}$ to $\ol{v}$ with intermediates in $W_2$ either (A) consists of a single fresh edge $(I_w, w)$ for some $w \in W_1 \sm J$, or (B) uses only original edges from $E$.}

The fresh edge $(I_w, w)$ has source $I_w$ (type-disjoint from $J \cup V \supseteq W_2$, hence $I_w \notin W_2$) and target $w \in W_1 \sm J$ (with $W_1 \cap W_2 = \emptyset$, hence $w \notin W_2$). We check where in the walk a fresh edge can appear:
\begin{itemize}
    \item \emph{Interior position.} If a fresh edge $(I_w, w)$ appears between two intermediate nodes (i.e.\ as $(a_i, a_{i+1})$ for some $1 \le i \le n - 1$), then $a_i = I_w$ or $a_{i+1} = w$ would force an intermediate node to lie outside $W_2$ (since neither $I_w$ nor $w$ lies in $W_2$). Contradiction.
    \item \emph{Last edge ($n \ge 2$).} If $(a_{n-1}, \ol{v}) = (I_w, w)$, then $a_{n-1} = I_w$, but $a_{n-1}$ is intermediate and so must lie in $W_2$; $I_w \notin W_2$. Contradiction.
    \item \emph{First edge ($n \ge 2$).} If $(\ul{v}, a_1) = (I_w, w)$, then $a_1 = w \in W_1 \sm J$; $a_1$ is intermediate and so must lie in $W_2$; $W_1 \cap W_2 = \emptyset$ gives $w \notin W_2$. Contradiction.
    \item \emph{First edge ($n = 1$, the whole walk).} The walk is the single edge $(\ul{v}, \ol{v}) = (I_w, w)$, with no intermediate-node constraint to satisfy; this case is admitted and produces case (A).
\end{itemize}
Hence either the walk is a single fresh edge (case (A)), or no fresh edge appears anywhere in the walk and every edge lies in $E$ (case (B)).

Case (A) contributes the pairs $\lC (I_w, w) \st w \in W_1 \sm J \rC$ to the LHS edge set; each such pair has $I_w$ in the input slot ($I_w \in J_{\lp G_{\doit(I_{W_1})} \rp^{\sm W_2}}$ by the input-node calculation above) and $w$ in the output slot ($w \in W_1 \sm J \ins V$, and $w \notin W_2$ by disjointness, so $w \in V \sm W_2 = V_{\lp G_{\doit(I_{W_1})} \rp^{\sm W_2}}$).

Case (B) contributes the pairs $(\ul{v}, \ol{v}) \in (J \cup (V \sm W_2)) \times (V \sm W_2)$ such that there is a directed walk in $G$ from $\ul{v}$ to $\ol{v}$ with intermediate nodes in $W_2$, i.e.\ exactly $E^{\sm W_2}$ (by def \ref{def:G_marginalization}, item~iii.); the source $\ul{v}$ must lie in $J \cup V$ because $E \ins (J \cup V) \times V$, and the absence of fresh edges in the walk excludes the intervention-symbol part of the carrier from the source. The typing $\ul{v} \notin W_2$ and $\ol{v} \notin W_2$ comes from $\ul{v}, \ol{v}$ lying in the marginalized output side, combined with the fact that any source in $V$ that contributed to a walk-through-$W_2$ pair survives the marginalization typing.

Together:
\[
    (E_{\doit(I_{W_1})})^{\sm W_2} \;=\; E^{\sm W_2} \;\cup\; \lC (I_w, w) \st w \in W_1 \sm J \rC.
\]

RHS: by def \ref{def:cdmg_intervention_nodes}, item~iii., applied to $G^{\sm W_2}$,
\[
    (E^{\sm W_2})_{\doit(I_{W_1})} \;=\; E^{\sm W_2} \;\cup\; \lC (I_w, w) \st w \in W_1 \sm J^{\sm W_2} \rC \;=\; E^{\sm W_2} \;\cup\; \lC (I_w, w) \st w \in W_1 \sm J \rC,
\]
using $J^{\sm W_2} = J$.

The LHS and RHS edge sets coincide. \checkmark

\medskip\noindent\emph{Bidirected edges.}
By def \ref{def:cdmg_intervention_nodes}, item~iv., $L_{G_{\doit(I_{W_1})}} = L$: adding intervention nodes does not touch the bidirected edge set. Hence on the LHS,
\[
    (L_{\doit(I_{W_1})})^{\sm W_2} \;=\; L^{\sm W_2}
\]
(the inner extension is a no-op on $L$, so the outer marginalization sees the original $L$). On the RHS, by def \ref{def:cdmg_intervention_nodes}, item~iv., applied to $G^{\sm W_2}$,
\[
    (L^{\sm W_2})_{\doit(I_{W_1})} \;=\; L^{\sm W_2}
\]
(the outer extension is a no-op on $L^{\sm W_2}$). Equal. \checkmark

\smallskip
Combining the four field-by-field equalities, $\lp G_{\doit(I_{W_1})} \rp^{\sm W_2} = \lp G^{\sm W_2} \rp_{\doit(I_{W_1})}$ as CDMGs. This establishes Part (ii).

\medskip\noindent\textbf{Part (iii) -- node-splitting.}
We verify that all four components of $\lp G_{\spl(W_1)} \rp^{\sm W_2}$ and $\lp G^{\sm W_2} \rp_{\spl(W_1)}$ agree, with $W_1 \ins V$, $W_2 \ins V$, and $W_1 \cap W_2 = \emptyset$. The construction of $G_{\spl(W_1)}$ is by def \ref{def:G_node-splitting}: tagged copies $W_1^0 := \lC w^0 \st w \in W_1 \rC$ and $W_1^1 := \lC w^1 \st w \in W_1 \rC$ of $W_1$ are introduced (type-disjoint from $J \cup V$ and from each other), with the shorthand $v^0 := v^1 := v$ for $v \in J \cup (V \sm W_1)$.

\smallskip\noindent\emph{Well-typedness.}
The inner $\spl(W_1)$ on the LHS is defined per def \ref{def:G_node-splitting} from $W_1 \ins V$. By items~i.\ and ii.\ of def \ref{def:G_node-splitting}, $G_{\spl(W_1)}$ has input nodes $J$ and output nodes $(V \sm W_1) \dcup W_1^0 \dcup W_1^1$. The outer $(\cdot)^{\sm W_2}$ requires $W_2$ to be a subset of the output-node set of $G_{\spl(W_1)}$; this holds because $W_2 \ins V$ and $W_1 \cap W_2 = \emptyset$ give $W_2 \ins V \sm W_1$, which embeds into $(V \sm W_1) \dcup W_1^0 \dcup W_1^1$ via the untagged-shorthand identification $v^0 = v^1 = v$ of def \ref{def:G_node-splitting} (for $v \in V \sm W_1$). Throughout the proof we identify $W_2$ with its untagged image in $V_{\spl(W_1)}$. Symmetrically on the RHS: $G^{\sm W_2}$ has output-node set $V \sm W_2$; the outer $\spl(W_1)$ requires $W_1 \ins V \sm W_2$, which holds because $W_1 \ins V$ and $W_1 \cap W_2 = \emptyset$.

\medskip\noindent\emph{Input nodes.}
Both sides have $J$: node-splitting preserves $J$ (def \ref{def:G_node-splitting}, item~i.\ gives $J_{G_{\spl(W_1)}} = J$), and marginalization preserves $J$ (def \ref{def:G_marginalization}, item~i.). Hence
\[
    J_{\lp G_{\spl(W_1)} \rp^{\sm W_2}} \;=\; J_{G_{\spl(W_1)}} \;=\; J \;=\; J^{\sm W_2} \;=\; J_{\lp G^{\sm W_2} \rp_{\spl(W_1)}}. \quad \checkmark
\]

\medskip\noindent\emph{Output nodes.}
LHS: by def \ref{def:G_marginalization}, item~ii., applied to $G_{\spl(W_1)}$,
\[
    V_{\lp G_{\spl(W_1)} \rp^{\sm W_2}} \;=\; V_{G_{\spl(W_1)}} \sm W_2 \;=\; \lp (V \sm W_1) \dcup W_1^0 \dcup W_1^1 \rp \sm W_2.
\]
Since $W_2 \ins V \sm W_1$ and $W_1^0, W_1^1$ are type-disjoint from $J \cup V$ (and so from $W_2 \ins V$), the removal $\sm W_2$ acts only on the $V \sm W_1$ piece:
\[
    \lp (V \sm W_1) \dcup W_1^0 \dcup W_1^1 \rp \sm W_2 \;=\; \lp (V \sm W_1) \sm W_2 \rp \dcup W_1^0 \dcup W_1^1 \;=\; \lp V \sm (W_1 \cup W_2) \rp \dcup W_1^0 \dcup W_1^1.
\]
RHS: by def \ref{def:G_node-splitting}, item~ii., applied to $G^{\sm W_2}$ (whose output-node set is $V \sm W_2$ with the tagged copies of $W_1$ taken inside $V \sm W_2$),
\[
    V_{\lp G^{\sm W_2} \rp_{\spl(W_1)}} \;=\; \lp (V \sm W_2) \sm W_1 \rp \dcup W_1^0 \dcup W_1^1 \;=\; \lp V \sm (W_1 \cup W_2) \rp \dcup W_1^0 \dcup W_1^1.
\]
Both sides coincide. \checkmark

\medskip\noindent\emph{Directed edges.}
Throughout this paragraph we work over the carrier of $G_{\spl(W_1)}$, namely $J \cup (V \sm W_1) \dcup W_1^0 \dcup W_1^1$. The subset $W_2 \ins V \sm W_1$ embeds into this carrier via the untagged shorthand of def \ref{def:G_node-splitting} (each $w \in W_2$ identifies with itself in $V \sm W_1$, distinct from every $w'^0$ and $w'^1$ for $w' \in W_1$).

By def \ref{def:G_node-splitting}, item~iii., $E_{\spl(W_1)}$ decomposes as the disjoint union of the \emph{lifted edges} $\lC (v_1^1, v_2^0) \st (v_1, v_2) \in E \rC$ and the \emph{transfer edges} $\lC (w^0, w^1) \st w \in W_1 \rC$.

Let $(x, y)$ in the appropriate slots of $(E_{\spl(W_1)})^{\sm W_2}$, i.e.\ $x \in J_{\lp G_{\spl(W_1)} \rp^{\sm W_2}} \cup V_{\lp G_{\spl(W_1)} \rp^{\sm W_2}} = J \cup (V \sm (W_1 \cup W_2)) \dcup W_1^0 \dcup W_1^1$ and $y \in V_{\lp G_{\spl(W_1)} \rp^{\sm W_2}}$.

LHS: $(x, y) \in (E_{\spl(W_1)})^{\sm W_2}$ iff (by def \ref{def:G_marginalization}, item~iii.) there is a directed walk
\[ x \tuh a_1 \tuh \cdots \tuh a_{n-1} \tuh y \]
in $G_{\spl(W_1)}$ with $n \ge 1$ and intermediate nodes $a_1, \dots, a_{n-1} \in W_2$.

\emph{Claim: any walk in $G_{\spl(W_1)}$ from $x$ to $y$ with intermediates in $W_2$ either (A) consists of a single transfer edge $(w^0, w^1)$ for some $w \in W_1$, or (B) uses only lifted edges.}

Each intermediate node $a_i \in W_2 \ins V \sm W_1$ is \emph{untagged} (in the def \ref{def:G_node-splitting} disjoint-union decomposition, every element of $V \sm W_1$ sits in the untagged piece, type-disjoint from $W_1^0$ and from $W_1^1$). We check where a transfer edge can appear:
\begin{itemize}
    \item \emph{Interior position ($n \ge 3$).} If $(a_i, a_{i+1}) = (w^0, w^1)$ for $1 \le i \le n - 2$, then $a_i = w^0 \in W_1^0$, but $a_i$ is intermediate and must lie in $W_2 \ins V \sm W_1$ (untagged). Type-disjointness of $W_1^0$ from $V \sm W_1$ gives a contradiction.
    \item \emph{Last edge ($n \ge 2$).} If $(a_{n-1}, y) = (w^0, w^1)$, then $a_{n-1} = w^0 \in W_1^0$, but $a_{n-1}$ is intermediate, contradiction (same type-disjointness).
    \item \emph{First edge ($n \ge 2$).} If $(x, a_1) = (w^0, w^1)$, then $a_1 = w^1 \in W_1^1$, but $a_1$ is intermediate, contradiction.
    \item \emph{First edge ($n = 1$, the whole walk).} The walk is the single transfer edge $(x, y) = (w^0, w^1)$, with no intermediate-node constraint to satisfy; this case is admitted and produces case (A).
\end{itemize}
Hence either the walk is a single transfer edge (case (A)), or no transfer edge appears anywhere in the walk and every edge is a lifted edge (case (B)).

\emph{Case (A) contribution.} The transfer edges $\lC (w^0, w^1) \st w \in W_1 \rC$ appear directly in $(E_{\spl(W_1)})^{\sm W_2}$: each $(w^0, w^1)$ has $w^0 \in W_1^0$ in the input/output side (specifically $w^0 \in V_{\spl(W_1)} \ins J_{\spl(W_1)} \cup V_{\spl(W_1)}$, and $w^0 \notin W_2$ by type-disjointness of $W_1^0$ from $V$) and $w^1 \in W_1^1 \ins V_{\spl(W_1)}$ with $w^1 \notin W_2$ by the same type-disjointness; both survive the marginalization output-side restriction $\sm W_2$.

\emph{Case (B) analysis.} Each lifted edge $(v_a^1, v_b^0)$ corresponds to an original edge $(v_a, v_b) \in E$, with $v_a \in J \cup V$ and $v_b \in V$ (by the typing of $E$). For a walk of lifted edges $x = x_0 \tuh x_1 \tuh \cdots \tuh x_n = y$ with $x_i$ intermediate ($1 \le i \le n - 1$) lying in $W_2 \ins V \sm W_1$:
\begin{itemize}
    \item Each intermediate $x_i$ is untagged: $x_i \in V \sm W_1$ and (by the untagged-shorthand convention of def \ref{def:G_node-splitting}) $x_i^0 = x_i^1 = x_i$.
    \item The first edge $(x_0, x_1) = (v_1^1, v_2^0)$ for some $(v_1, v_2) \in E$, hence $v_2 = x_1$ untagged forces $v_2 \in V \sm W_1$ and $v_2^0 = v_2 = x_1$.
    \item Consecutive lifted edges $(x_{i-1}, x_i) = (u_a^1, u_b^0)$ and $(x_i, x_{i+1}) = (u_a'^1, u_b'^0)$ with $u_b = x_i = u_a'$ (untagged) match coherently because $u_b^0 = u_b = u_a' = u_a'^1$ at the untagged intermediate $x_i$. The underlying edges $(u_a, u_b), (u_a', u_b') \in E$ thus chain into a walk in $G$.
    \item The last edge $(x_{n-1}, x_n) = (v_n^1, v_{n+1}^0)$ for some $(v_n, v_{n+1}) \in E$, with $v_n = x_{n-1}$ untagged (for $n \ge 2$; if $n = 1$, $x_{n-1} = x_0$ is the source endpoint and need not be untagged).
\end{itemize}

The endpoints $x = x_0$ and $y = x_n$ may be tagged or untagged. From the form of the first/last lifted edges:
\begin{itemize}
    \item $x = v_1^1$: if $v_1 \in W_1$, then $x = v_1^1 \in W_1^1$; if $v_1 \in J \cup (V \sm W_1)$, then $x = v_1$ untagged.
    \item $y = v_{n+1}^0$: if $v_{n+1} \in W_1$, then $y = v_{n+1}^0 \in W_1^0$; if $v_{n+1} \in V \sm W_1$, then $y = v_{n+1}$ untagged.
\end{itemize}

The typing of $(x, y)$ in $(E_{\spl(W_1)})^{\sm W_2}$ forces, in the untagged case, $v_1 \notin W_2$ (so $v_1 \in J \cup (V \sm (W_1 \cup W_2))$) and $v_{n+1} \notin W_2$ (so $v_{n+1} \in V \sm (W_1 \cup W_2)$). In the tagged case, $v_1 \in W_1$ automatically gives $v_1 \notin W_2$ by $W_1 \cap W_2 = \emptyset$; similarly $v_{n+1} \in W_1$ gives $v_{n+1} \notin W_2$.

Hence in all sub-cases, $v_1 \in J \cup (V \sm W_2)$ and $v_{n+1} \in V \sm W_2$, and the walk $v_1 \tuh v_2 \tuh \cdots \tuh v_{n+1}$ in $G$ has intermediate nodes $v_2, \dots, v_n$ (for $n \ge 2$; empty for $n = 1$), each equal to the corresponding $x_i$ untagged, hence in $W_2$. By def \ref{def:G_marginalization}, item~iii., $(v_1, v_{n+1}) \in E^{\sm W_2}$.

Together (cases (A) and (B)):
\[
    (E_{\spl(W_1)})^{\sm W_2} \;=\; \lC (v_1^1, v_2^0) \st (v_1, v_2) \in E^{\sm W_2} \rC \;\cup\; \lC (w^0, w^1) \st w \in W_1 \rC.
\]

RHS: by def \ref{def:G_node-splitting}, item~iii., applied to $G^{\sm W_2}$ (whose directed-edge set is $E^{\sm W_2}$),
\[
    (E^{\sm W_2})_{\spl(W_1)} \;=\; \lC (v_1^1, v_2^0) \st (v_1, v_2) \in E^{\sm W_2} \rC \;\cup\; \lC (w^0, w^1) \st w \in W_1 \rC.
\]
(The transfer edges $\lC (w^0, w^1) \st w \in W_1 \rC$ are preserved verbatim across both orderings: on the LHS they survive marginalization as direct edges via case (A) of the claim above; on the RHS they are introduced by def \ref{def:G_node-splitting}, item~iii.'s second set-builder applied to the marginalized graph.)

The LHS and RHS edge sets coincide. \checkmark

\medskip\noindent\emph{Bidirected edges.}
By def \ref{def:G_node-splitting}, item~iv., the bidirected-edge set of $G_{\spl(W_1)}$ lives on the $^0$-tagged side only:
\[
    L_{\spl(W_1)} \;=\; \lC (v_1^0, v_2^0) \st (v_1, v_2) \in L \rC.
\]
No bidirected edge of $G_{\spl(W_1)}$ touches $W_1^1$: every endpoint of a bidirected edge carries the $^0$-superscript (and is hence either untagged via the shorthand $v^0 = v$ for $v \in V \sm W_1$, or tagged into $W_1^0$ for $v \in W_1$).

Let $u, v$ in the slots of $(L_{\spl(W_1)})^{\sm W_2}$: $u, v \in V_{\lp G_{\spl(W_1)} \rp^{\sm W_2}} = (V \sm (W_1 \cup W_2)) \dcup W_1^0 \dcup W_1^1$ with $u \neq v$ (def \ref{def:G_marginalization}, item~iv.).

LHS: $u \huh v \in (L_{\spl(W_1)})^{\sm W_2}$ iff (by def \ref{def:G_marginalization}, item~iv.) there is a bifurcation in $G_{\spl(W_1)}$ between $u$ and $v$ with intermediate nodes in $W_2$.

\emph{Excluding $W_1^1$-endpoints.} Suppose $u \in W_1^1$, say $u = w^1$ for some $w \in W_1$. We show no bifurcation in $G_{\spl(W_1)}$ ending at $u$ with intermediates in $W_2$ can exist. Tracing the def \ref{def:G_marginalization}, item~iv., bifurcation structure $(u = w_0, w_1, \ldots, w_n = v)$ with hinge at index $k$:
\begin{itemize}
    \item If $k \ge 2$, the first left-arm edge $(w_1, w_0) = (w_1, u) \in E_{\spl(W_1)}$. A lifted edge $(v_a^1, v_b^0)$ here gives $u = v_b^0 \in (V \sm W_1) \cup W_1^0$, never in $W_1^1$ -- contradiction. A transfer edge gives source $w_1 = w'^0 \in W_1^0$, which is intermediate (if $n \ge 2$) and must lie in $W_2 \ins V \sm W_1$, type-disjoint from $W_1^0$ -- contradiction.
    \item If $k = 1$, the hinge is at index 1. The directed-hinge alternative $(w_1, w_0) = (w_1, u) \in E_{\spl(W_1)}$ runs into the same lifted/transfer-edge contradictions as the $k \ge 2$ case. The bidirected-hinge alternative $(w_0, w_1) = (u, w_1) \in L_{\spl(W_1)}$ gives $u = v_1^0$ for some $(v_1, v_2) \in L$, placing $u \in (V \sm W_1) \cup W_1^0$, never in $W_1^1$ -- contradiction.
\end{itemize}
Hence no bifurcation in $G_{\spl(W_1)}$ between $u \in W_1^1$ and any $v$ with intermediates in $W_2$ exists. By symmetry (def \ref{def:G_marginalization}, item~iv.'s bifurcation is symmetric in $u$ and $v$ via reverse traversal -- equivalently, the bifurcation predicate $\Phi_L$ is symmetric in its two arguments per the symmetry axiom of def \ref{def-cdmg} transported through def \ref{def:G_marginalization}, item~iv.), the same conclusion holds with $v \in W_1^1$. So the effective range of bidirected edges in $(L_{\spl(W_1)})^{\sm W_2}$ has both endpoints in $(V \sm (W_1 \cup W_2)) \cup W_1^0$.

\emph{Correspondence with bifurcations in $G$ through $W_2$.} Suppose $u, v \in (V \sm (W_1 \cup W_2)) \cup W_1^0$. By the same analysis as for directed edges (transfer edges cannot appear because their endpoints are tagged in $W_1^0 \cup W_1^1$, type-disjoint from $W_2 \ins V \sm W_1$ for intermediates), every directed-arm edge of the bifurcation is a lifted edge $(v_a^1, v_b^0)$ from some $(v_a, v_b) \in E$, and the bidirected-hinge edge (if any) is a lifted bidirected edge $(v_1^0, v_2^0)$ from some $(v_1, v_2) \in L$.

For each intermediate $w_i$ in $W_2 \ins V \sm W_1$ (untagged): if $w_i$ appears as the target of one arm edge with $^0$-tag and as the source of the next arm edge with $^1$-tag, the untagged-shorthand convention $w_i^0 = w_i^1 = w_i$ makes the two tags coherent. Tracking the $^0$/$^1$ tags through the bifurcation skeleton of def \ref{def:G_marginalization}, item~iv.:
\begin{itemize}
    \item Left-arm edges $(w_i, w_{i-1}) \in E_{\spl(W_1)}$ for $1 \le i \le k - 1$ lift to $(w_i^1, w_{i-1}^0) \in E_{\spl(W_1)}$ from $(w_i, w_{i-1}) \in E$, putting $^0$ on $w_{i-1}$ and $^1$ on $w_i$.
    \item Right-arm edges $(w_i, w_{i+1}) \in E_{\spl(W_1)}$ for $k \le i \le n - 1$ lift similarly, putting $^0$ on $w_{i+1}$ and $^1$ on $w_i$.
    \item The hinge (directed alternative $(w_k, w_{k-1}) \in E_{\spl(W_1)}$ or bidirected alternative $(w_{k-1}, w_k) \in L_{\spl(W_1)}$) lifts respectively to $(w_k^1, w_{k-1}^0)$ or $(w_{k-1}^0, w_k^0)$.
\end{itemize}
Hence both endpoints $u = w_0$ and $v = w_n$ acquire the $^0$-superscript (from being the target of the first left-arm edge, the right endpoint of a bidirected hinge at $k = 1$, the target of the right-arm last edge, or the right endpoint of a bidirected hinge at $k = n$). Intermediate nodes $w_i$ for $1 \le i \le n - 1$ acquire \emph{both} $^0$ and $^1$ superscripts (consistent only if $w_i$ is untagged, which holds because $w_i \in W_2 \ins V \sm W_1$). For the hinge node $w_k$: if the hinge is directed, $w_k$ appears as the source of the hinge edge with $^1$-tag and (if $k \le n - 1$) as the source of the right-arm edge with $^1$-tag -- coherent; if the hinge is bidirected, $w_k$ appears with $^0$-tag from the hinge edge and (if $k \le n - 1$) with $^1$-tag from the right-arm edge -- coherent only if $w_k$ is untagged, which holds at intermediate positions.

The above lifting argument runs both ways:
\begin{itemize}
    \item Every bifurcation in $G_{\spl(W_1)}$ between $u, v \in (V \sm (W_1 \cup W_2)) \cup W_1^0$ with intermediates in $W_2$ corresponds to a bifurcation in $G$ between $u', v' \in V \sm W_2$ with intermediates in $W_2$, where $u'$ is the underlying node of $u$ (so $u = u'^0$) and $v'$ is the underlying node of $v$ (so $v = v'^0$). Here $u' \neq v'$ since $u \neq v$ and the map $v \mapsto v^0$ is injective on $V$ (def \ref{def:G_node-splitting}, item~ii., tagged construction).
    \item Conversely, every bifurcation in $G$ between $u', v' \in V \sm W_2$ with $u' \neq v'$ and intermediates in $W_2$ lifts to a bifurcation in $G_{\spl(W_1)}$ between $u'^0, v'^0$ with intermediates in $W_2 \ins V \sm W_1$ (each directed-arm edge of $E$ lifts to a lifted directed edge of $E_{\spl(W_1)}$, the bidirected-hinge edge of $L$ lifts to a lifted bidirected edge of $L_{\spl(W_1)}$, and intermediates retain their untagged identity via the shorthand $w_i^0 = w_i^1 = w_i$).
\end{itemize}

By def \ref{def:G_marginalization}, item~iv., applied to $G$ (with $u', v' \in V \sm W_2$, $u' \neq v'$, bifurcation through $W_2$), the bifurcation-through-$W_2$ condition on $G$ is exactly the predicate defining $L^{\sm W_2}$.

Hence
\[
    (L_{\spl(W_1)})^{\sm W_2} \;=\; \lC (v_1^0, v_2^0) \st (v_1, v_2) \in L^{\sm W_2} \rC.
\]
(Bidirected edges with at least one endpoint in $W_1^1$ are absent on both sides by the analysis above; the irreflexivity constraint $u \neq v$ on the LHS matches the irreflexivity of $L^{\sm W_2}$ via the injectivity of $v \mapsto v^0$.)

RHS: by def \ref{def:G_node-splitting}, item~iv., applied to $G^{\sm W_2}$ (whose bidirected-edge set is $L^{\sm W_2}$),
\[
    (L^{\sm W_2})_{\spl(W_1)} \;=\; \lC (v_1^0, v_2^0) \st (v_1, v_2) \in L^{\sm W_2} \rC.
\]
Equal. \checkmark

\smallskip
Combining the four field-by-field equalities, $\lp G_{\spl(W_1)} \rp^{\sm W_2} = \lp G^{\sm W_2} \rp_{\spl(W_1)}$ as CDMGs. This establishes Part (iii).

\medskip
Parts (i), (ii), and (iii) jointly establish the three displayed equalities of the lemma.
\end{proof}

\end{document}
